% In this file you should put the actual content of the blueprint.
% It will be used both by the web and the print version.
% It should *not* include the \begin{document}
%
% If you want to split the blueprint content into several files then
% the current file can be a simple sequence of \input. Otherwise It
% can start with a \section or \chapter for instance.

\chapter{Spectral triples}
\label{cha:spectral-triples}

This chapter records the core operator predicates underlying odd and even spectral triples
and the first consequences of their axioms.  The predicates deliberately separate the
self-adjointness, domain, commutator, and grading axioms from compactness of the resolvent;
the latter is added by Definition~\ref{def:compact-resolvent-spectral-triple}.  This split is
an implementation convenience, not the standard terminology in which compact resolvent is
already part of the definition of a spectral triple.

\section{Odd spectral triples}
\label{sec:odd-spectral-triples}

\begin{definition}[Odd spectral-triple core]
  \label{def:odd-spectral-triple}
  \lean{IsOddSpectralTriple}
  \leanok
  Let $\Hil$ be a Hilbert space over $\Bbbk \in \{\mathbb{R}, \mathbb{C}\}$, let $A$ be a
  $*$-algebra over $\Bbbk$, and let $\pi \colon A \to \mathcal{B}(\Hil)$ be a
  $*$-algebra homomorphism. A possibly unbounded linear operator
  $D \colon \dom{D} \to \Hil$, with $\dom{D} \subseteq \Hil$ a linear subspace,
  together with the data $(A, \Hil, \pi, D)$, satisfies the \emph{odd spectral-triple core}
  predicate if:
  \begin{itemize}
    \item $D$ is self-adjoint (in particular, $\dom{D}$ is dense and $D$ is closed);
    \item for every $a \in A$, $\pi(a)$ maps $\dom{D}$ into $\dom{D}$;
    \item for every $a \in A$, the commutator $[D, \pi(a)]$, viewed as an operator on
      $\dom{D}$, is bounded on the closed unit ball of $\dom{D}$.
  \end{itemize}
  No compact-resolvent or Schatten-summability condition is included in this predicate.
\end{definition}

\begin{lemma}[Density of the domain of the Dirac operator]
  \label{lem:dense-domain-dirac}
  \lean{IsOddSpectralTriple.dense_domain_dirac}
  \leanok
  \uses{def:odd-spectral-triple}
  If $(A, \Hil, \pi, D)$ satisfies the odd core predicate, then $\dom{D}$ is dense in $\Hil$.
\end{lemma}

\begin{lemma}[Closedness of the Dirac operator]
  \label{lem:isClosed-dirac}
  \lean{IsOddSpectralTriple.isClosed_dirac}
  \leanok
  \uses{def:odd-spectral-triple}
  If $(A, \Hil, \pi, D)$ satisfies the odd core predicate, then $D$ is a closed operator.
\end{lemma}

\begin{definition}[The algebraic commutator on the Dirac domain]
  \label{def:commutator-on-domain}
  \lean{IsOddSpectralTriple.domainRepresentation,
        IsOddSpectralTriple.commutatorOnDomain}
  \leanok
  \uses{def:odd-spectral-triple}
  Domain invariance restricts $\pi(a)$ to a linear endomorphism of $\dom{D}$.  The algebraic
  commutator on that domain is the linear map
  \[
    C_a\xi = D\pi(a)\xi-\pi(a)D\xi = [D,\pi(a)]\xi.
  \]
\end{definition}

\begin{lemma}[Unit-ball and homogeneous commutator bounds]
  \label{lem:exists-comm-bound}
  \lean{IsOddSpectralTriple.exists_comm_bound,
        IsOddSpectralTriple.exists_commutator_bound}
  \leanok
  \uses{def:odd-spectral-triple, def:commutator-on-domain}
  If $(A, \Hil, \pi, D)$ satisfies the odd core predicate, then for every $a \in A$ there is
  a real constant $C \geq 0$ such that
  \[
    \|C_a\xi\| \leq C\|\xi\|
  \]
  for every $\xi \in \dom{D}$.  This homogeneous estimate is obtained by scaling the
  unit-ball bound in the core axiom.
\end{lemma}

\begin{definition}[The bounded commutator]
  \label{def:bounded-commutator}
  \lean{IsOddSpectralTriple.boundedCommutator,
        IsOddSpectralTriple.boundedCommutator_apply,
        IsOddSpectralTriple.boundedCommutator_unique}
  \leanok
  \uses{lem:dense-domain-dirac, def:commutator-on-domain, lem:exists-comm-bound}
  The homogeneous estimate extends $C_a$ uniquely from the dense domain $\dom{D}$ to a
  continuous linear map $\overline C_a\colon\Hil\to\Hil$.  On the original domain it satisfies
  \[
    \overline C_a\xi=D\pi(a)\xi-\pi(a)D\xi.
  \]
\end{definition}

\section{Even spectral triples}
\label{sec:even-spectral-triples}

\begin{definition}[Even spectral-triple core]
  \label{def:even-spectral-triple}
  \lean{IsEvenSpectralTriple}
  \leanok
  \uses{def:odd-spectral-triple}
  The \emph{even spectral-triple core} predicate for $(A, \Hil, \pi, D, \gamma)$ consists of
  the odd core predicate for $(A, \Hil, \pi, D)$ together with a bounded operator
  $\gamma \in \mathcal{B}(\Hil)$ (the \emph{grading operator}) such that:
  \begin{itemize}
    \item $\gamma$ is self-adjoint: $\gamma^* = \gamma$;
    \item $\gamma$ is an involution: $\gamma^2 = 1$;
    \item $\gamma$ commutes with $\pi(a)$ for every $a \in A$;
    \item $\gamma$ maps $\dom{D}$ into $\dom{D}$;
    \item $D \gamma \xi = -\gamma D \xi$ for every $\xi \in \dom{D}$.
  \end{itemize}
\end{definition}

\section{Resolvents and compact-resolvent spectral triples}
\label{sec:resolvents}

This section introduces the resolvent set and resolvent of the (partially defined) Dirac
operator $D$, and adds the compact-resolvent condition to the core predicate.  Compact
resolvent alone is weaker than finite or $p$-summability, which would require a Schatten- or
trace-ideal estimate and is not formalized here.

\begin{definition}[Resolvent set]
  \label{def:resolvent-set}
  \lean{LinearPMap.resolventSet}
  \leanok
  Let $D \colon \dom{D} \to \Hil$ be a (partially defined) linear operator, where
  $\dom{D} \subseteq \Hil$ is a linear subspace. For $z \in \Bbbk$, write $z - D$ for the
  partially defined linear map $\xi \mapsto z\xi - D\xi$ with domain $\dom{D}$. The
  \emph{resolvent set} $\rho(D) \subseteq \Bbbk$ of $D$ is the set of $z \in \Bbbk$ for which
  $z - D \colon \dom{D} \to \Hil$ is bijective.
\end{definition}

\begin{definition}[Resolvent]
  \label{def:resolvent}
  \lean{LinearPMap.resolvent}
  \leanok
  \uses{def:resolvent-set}
  For $z \in \rho(D)$, the \emph{resolvent} $R(z, D) \colon \Hil \to \Hil$ is the
  (everywhere-defined, linear) inverse of $z - D$.
\end{definition}

\begin{lemma}[Range of the resolvent]
  \label{lem:range-resolvent}
  \lean{LinearPMap.range_resolvent}
  \leanok
  \uses{def:resolvent}
  For $z \in \rho(D)$, the range of $R(z, D)$ equals $\dom{D}$.
\end{lemma}

\begin{lemma}[Norm bound for self-adjoint operators]
  \label{lem:norm-resolvent-apply-ge}
  \lean{IsSelfAdjoint.norm_resolvent_apply_ge}
  \leanok
  If $D$ is self-adjoint and $z \in \Bbbk$, then for every $\xi \in \dom{D}$,
  \[
    |\operatorname{Im}(z)| \cdot \|\xi\| \leq \|z\xi - D\xi\|.
  \]
\end{lemma}

\begin{lemma}[Injectivity of $z - D$ off the real axis]
  \label{lem:injective-resolvent-apply}
  \lean{IsSelfAdjoint.injective_resolvent_apply}
  \leanok
  \uses{lem:norm-resolvent-apply-ge}
  If $D$ is self-adjoint and $\operatorname{Im}(z) \neq 0$, then $z - D \colon \dom{D} \to \Hil$
  is injective.
\end{lemma}

\begin{theorem}[Basic criterion of self-adjointness]
  \label{thm:basic-criterion-self-adjoint}
  \lean{IsSelfAdjoint.mem_resolventSet}
  \leanok
  \uses{lem:injective-resolvent-apply}
  If $D$ is self-adjoint and $\operatorname{Im}(z) \neq 0$, then $z - D$ is in fact bijective,
  i.e.\ $z \in \rho(D)$.

  This strengthens Lemma~\ref{lem:injective-resolvent-apply} from injectivity to bijectivity:
  the range of $z - D$ is closed (the bounded-below estimate makes preimage sequences Cauchy,
  and the graph of $D$ is closed) and dense (its orthogonal complement is trivial, via the
  adjoint), hence all of $\Hil$. In particular $i \in \rho(D)$, so the resolvent-set hypothesis
  needed by the compact-resolvent predicate holds automatically at $z = i$ over $\mathbb C$.
\end{theorem}

\begin{definition}[Compact-resolvent spectral triple]
  \label{def:compact-resolvent-spectral-triple}
  \lean{IsCompactResolventSpectralTriple}
  \leanok
  \uses{def:odd-spectral-triple, def:resolvent-set, def:resolvent,
        thm:basic-criterion-self-adjoint}
  An odd spectral-triple core $(A, \Hil, \pi, D)$ has \emph{compact resolvent at $i$} if
  $i \in \rho(D)$ and the resolvent $R(i,D)$ is a compact operator on $\Hil$.  The
  resolvent-membership proof is an explicit field: over $\mathbb C$ it follows from
  Theorem~\ref{thm:basic-criterion-self-adjoint}, whereas for $\Bbbk=\mathbb R$ the value
  \texttt{RCLike.I} is $0$ and membership must be supplied separately.
\end{definition}

\chapter{The Fredholm index pairing}
\label{cha:index-pairing}

This chapter develops foundations for the Fredholm index pairing.  The formalized
\emph{graded-kernel invariant} (Section~\ref{sec:graded-kernel-index}) is the difference of
the two graded kernel dimensions.  The still-open chiral bridge must construct $D^+$, prove it
Fredholm, and identify its Fredholm index with that invariant.  The general pairing against an
arbitrary $K_0$-class $[p]$ is likewise not yet formalized.

\section{Fredholm operators}
\label{sec:fredholm}

This section records a minimal, self-contained notion of Fredholm linear map, scoped to
exactly what the index pairing needs. The Mathlib version pinned by this repository does not
provide a finalized general theory of Fredholm operators; the definitions here are
deliberately narrower — stated for plain linear maps rather than continuous linear maps, so
that they apply uniformly to bounded operators and to the restriction of an unbounded
operator to its domain, as needed for the chiral Dirac operator $D^{+}$ — and are expected to
be replaced by an \texttt{import} once the upstream theory lands.

\begin{definition}[Fredholm linear map]
  \label{def:isFredholm}
  \lean{SpectralTriples.Fredholm.IsFredholm}
  \leanok
  A linear map $f \colon E \to F$ is \emph{Fredholm} if $\ker f$ is finite-dimensional, the
  range of $f$ is closed, and the range of $f$ has finite codimension (the cokernel
  $F / \operatorname{ran} f$ is finite-dimensional).
\end{definition}

\begin{definition}[Fredholm index]
  \label{def:fredholm-index}
  \lean{SpectralTriples.Fredholm.index}
  \leanok
  \uses{def:isFredholm}
  For a Fredholm linear map $f \colon E \to F$, its \emph{Fredholm index} is the integer
  \[
    \operatorname{ind}(f) = \dim(\ker f) - \dim(F / \operatorname{ran} f).
  \]
  The Lean definition is total on all linear maps; its intended dimension-counting
  interpretation is used together with an \texttt{IsFredholm} witness.
\end{definition}

\begin{lemma}[Index of a bijection]
  \label{lem:fredholm-index-of-bijective}
  \lean{SpectralTriples.Fredholm.isFredholm_of_bijective,
        SpectralTriples.Fredholm.index_of_bijective}
  \leanok
  \uses{def:isFredholm, def:fredholm-index}
  A bijective linear map with closed range is Fredholm, of index $0$.
\end{lemma}

\subsection{Compact perturbations of the identity}
\label{subsec:compact-perturbations}

This subsection proves the structural part of the Hilbert-space case of the classical
Riesz--Schauder theorem: $1 - K$ is Fredholm whenever $K$ is compact (finite kernel, closed
range, finite-dimensional cokernel). The further classical fact that the index is $0$ is not
proved here. None of this is yet in Mathlib (only the weaker spectral dichotomy for compact
operators is), so it is built here from scratch, as the analytic ingredient needed to
eventually show the chiral Dirac operator $D^{+}$ is Fredholm.

\begin{theorem}[Compact operators are approximable by finite rank]
  \label{thm:exists-finiteRank-norm-sub-lt}
  \lean{IsCompactOperator.exists_finiteRank_norm_sub_lt}
  \leanok
  A compact operator $K$ on a Hilbert space is approximable, in operator norm, by finite-rank
  operators: for every $\varepsilon > 0$ there is a finite-rank $F$ with
  $\|K - F\| \leq \varepsilon$. The approximant is built by composing $K$ with the orthogonal
  projection onto the span of a finite $\varepsilon$-net of the (relatively compact) image of
  the closed unit ball.
\end{theorem}

\begin{lemma}[Finite-rank operators have finite-rank adjoints]
  \label{lem:finiteDimensional-range-adjoint}
  \lean{ContinuousLinearMap.finiteDimensional_range_adjoint}
  \leanok
  If a continuous linear map on an inner product space has finite-dimensional range, so does
  its adjoint.
\end{lemma}

\begin{theorem}[The adjoint of a compact operator is compact]
  \label{thm:isCompactOperator-adjoint}
  \lean{IsCompactOperator.adjoint}
  \leanok
  \uses{thm:exists-finiteRank-norm-sub-lt, lem:finiteDimensional-range-adjoint}
  If $K$ is a compact operator on a Hilbert space, then its adjoint $K^{\dagger}$ is also
  compact: the adjoints of finite-rank approximants of $K$
  (Theorem~\ref{thm:exists-finiteRank-norm-sub-lt}) are themselves finite-rank
  (Lemma~\ref{lem:finiteDimensional-range-adjoint}), hence compact, and converge to
  $K^{\dagger}$ in norm since the adjoint operation is norm-preserving.
\end{theorem}

\begin{theorem}[Compact perturbations of the identity are Fredholm]
  \label{thm:isFredholm-one-sub}
  \lean{SpectralTriples.Fredholm.isFredholm_one_sub}
  \leanok
  \uses{def:isFredholm, thm:isCompactOperator-adjoint}
  The structural part of the classical Riesz--Schauder theorem: if $K$ is compact on a
  Hilbert space, then $1 - K$ is Fredholm. The kernel is finite-dimensional (it is the
  eigenspace of $K$ at the eigenvalue $1$); the range is closed, via a bounded-below estimate
  on the orthogonal complement of the kernel established by a contradiction argument using
  compactness of $K$; and the cokernel is finite-dimensional, via the adjoint identity
  $(1-K)^{\dagger} = 1 - K^{\dagger}$ (Theorem~\ref{thm:isCompactOperator-adjoint}) together
  with $(\operatorname{ran}(1-K))^{\perp} = \ker((1-K)^{\dagger})$. (The further classical fact
  that the index is $0$ is not proved here.)
\end{theorem}

\section{The graded-kernel invariant}
\label{sec:graded-kernel-index}

\begin{definition}[Kernel of the Dirac operator]
  \label{def:Dkernel}
  \lean{SpectralTriples.Dkernel}
  \leanok
  For a (partially defined) linear operator $D \colon \dom{D} \to \Hil$, the \emph{kernel}
  $\ker D \subseteq \Hil$ is the subspace of vectors $\xi \in \dom{D}$ with $D\xi = 0$.
\end{definition}

\begin{lemma}[The grading preserves the kernel]
  \label{lem:grading-mem-Dkernel}
  \lean{SpectralTriples.grading_mem_Dkernel}
  \leanok
  \uses{def:even-spectral-triple, def:Dkernel}
  If $(A, \Hil, \pi, D, \gamma)$ satisfies the even spectral-triple core predicate and
  $\xi \in \ker D$, then $\gamma \xi \in \ker D$.  This theorem proves preservation of the
  kernel; a separate formal theorem supplies the eigenspace decomposition of the ambient
  grading.
\end{lemma}

\begin{definition}[Graded-kernel invariant]
  \label{def:graded-kernel-index}
  \lean{SpectralTriples.gradedKernelIndex}
  \leanok
  \uses{lem:grading-mem-Dkernel}
  The \emph{graded-kernel invariant} of an even Dirac datum $(D, \gamma)$ is the integer
  \[
    \operatorname{gkIndex}(D, \gamma) \;=\; \dim (\ker D)^{+} - \dim (\ker D)^{-}.
  \]
  The corresponding invariant of an even core predicate is
  \lean{IsEvenSpectralTriple.gradedKernelIndex}\leanok.  No chiral operator or equality with
  a Fredholm index is asserted by this definition.
\end{definition}

\begin{theorem}[Finite-dimensionality of the kernel]
  \label{thm:finiteDimensional-Dkernel}
  \lean{SpectralTriples.finiteDimensional_Dkernel}
  \leanok
  \uses{def:compact-resolvent-spectral-triple, def:Dkernel}
  If $D$ is self-adjoint, its resolvent at $i$ is compact, and
  $\operatorname{Im}(i)\neq 0$ (in particular over $\mathbb C$), then $\ker D$ is
  finite-dimensional.
  Consequently $\operatorname{gkIndex}(D, \gamma)$
  (Definition~\ref{def:graded-kernel-index}) is a difference of genuine finite dimensions:
  for $\xi \in \ker D$, the resolvent $R(i, D)$
  satisfies $R(i,D)(i\xi) = \xi$, so $\ker D$ embeds into the $i^{-1}$-eigenspace of the
  compact operator $R(i, D)$, which is finite-dimensional by Riesz theory.
\end{theorem}

\begin{theorem}[Chiral Fredholm bridge (planned)]
  \label{thm:chiral-fredholm-bridge}
  \uses{def:even-spectral-triple, def:compact-resolvent-spectral-triple,
        def:graded-kernel-index, def:isFredholm, def:fredholm-index,
        thm:finiteDimensional-Dkernel}
  For an even spectral-triple core with compact resolvent, construct the chiral restriction
  $D^+\colon \dom(D)\cap\Hil^+\to\Hil^-$, prove that it is Fredholm, and prove
  \[
    \operatorname{ind}(D^+) = \operatorname{gkIndex}(D,\gamma).
  \]
  This is an open dependency: finite-dimensionality of $\ker D$ alone does not prove closed
  range or finite-dimensional cokernel for the chiral restriction.
\end{theorem}

\section{The general index pairing}
\label{sec:general-index-pairing}

This section is not yet formalized; it corresponds to Phase~2 of the project plan.  It builds
on the planned chiral bridge and treats an arbitrary self-adjoint idempotent representing a
class in $K_0(A)$.

\begin{definition}[Index pairing]
  \label{def:index-pairing}
  \uses{def:even-spectral-triple, def:compact-resolvent-spectral-triple,
        def:fredholm-index}
  Let $(A, \Hil, \pi, D, \gamma)$ satisfy the even core and compact-resolvent predicates, and let
  $p \in M_n(A)$ be a self-adjoint idempotent (a projection representing a class
  in $K_0(A)$). The \emph{index pairing} of $p$ with $(A, \Hil, \pi, D, \gamma)$ is
  the Fredholm index
  \[
    \langle [p], (\Hil, \pi, D, \gamma) \rangle
      \;=\; \operatorname{ind}\big( \pi(p) D^{+} \pi(p) \big),
  \]
  where $D^{+}$ denotes the restriction of $D$ to the $(+1)$-eigenspace of $\gamma$.
\end{definition}

\begin{theorem}[The index pairing is well defined]
  \label{thm:index-pairing-well-defined}
  \uses{def:index-pairing, thm:chiral-fredholm-bridge, def:isFredholm}
  In the situation of Definition~\ref{def:index-pairing}, the operator
  $\pi(p) D^{+} \pi(p)$, restricted to $\pi(p)\Hil \cap \dom{D}$, is Fredholm, so that
  $\langle [p], (\Hil, \pi, D, \gamma) \rangle \in \mathbb{Z}$ is well defined.
\end{theorem}

\chapter{Examples}
\label{cha:examples}

This chapter records concrete spectral triples built directly on the Fourier side, bypassing
spin-geometry and Sobolev-space machinery from the literature, together with standalone
Fredholm-index examples that illustrate Section~\ref{sec:fredholm} directly. The circle and
torus examples share the same diagonal-operator technology (Section~\ref{sec:diagonal-operators}):
the circle example (Section~\ref{sec:circle}) is a complete odd core with compact resolvent;
the two-torus example (Section~\ref{sec:torus}) is a complete even core with compact resolvent
and vanishing graded-kernel invariant. Sections~\ref{sec:shift} and~\ref{sec:magnetic-dirac}
give two bounded Fredholm examples with indices $-1$ and $k$. Sections~\ref{sec:theta-sections}
and~\ref{sec:fourier-holomorphic} give the purely function-theoretic count
$\dim H^0(L_k)=k$ and the vanishing of the negative-degree holomorphic-section space.  Their
interpretation as the kernel and cokernel of a geometric $L^2$ Dirac operator remains part of
the planned operator bridge. Section~\ref{sec:hermite-l2} supplies a normalized orthonormal
Hermite family; completeness is the next missing analytic node.

\section{Block-diagonal operators on \texorpdfstring{$\ell^2$}{l2}}
\label{sec:diagonal-operators}

Given a family of finite-dimensional Hilbert spaces $(G_i)_{i \in \alpha}$ and a uniformly
bounded family of block operators $T_i \colon G_i \to G_i$, this section builds the
associated block-diagonal operator on $\ell^2(\alpha; G)$ and proves the compactness
criterion used to show the Dirac resolvents in Sections~\ref{sec:circle}
and~\ref{sec:torus} are compact.

\begin{definition}[Block-diagonal operator]
  \label{def:diagL}
  \lean{lpDiag.diagL}
  \leanok
  Given a uniformly bounded family of operators $T_i \colon G_i \to G_i$ with $\|T_i\| \le C$
  for all $i$, the \emph{block-diagonal operator} $\operatorname{diagL}(T) \colon
  \ell^2(\alpha; G) \to \ell^2(\alpha; G)$ is defined by $(\operatorname{diagL}(T) a)_i =
  T_i(a_i)$.
\end{definition}

\begin{lemma}[Operator norm bound]
  \label{lem:norm-diagL-le}
  \lean{lpDiag.norm_diagL_le}
  \leanok
  \uses{def:diagL}
  $\|\operatorname{diagL}(T)\| \le C$ whenever $\|T_i\| \le C$ for all $i$.
\end{lemma}

\begin{theorem}[Compactness via finite-rank truncation]
  \label{thm:isCompactOperator-diagL-of-support-finite}
  \lean{lpDiag.isCompactOperator_diagL_of_support_finite}
  \leanok
  \uses{def:diagL}
  If only finitely many blocks $T_i$ are nonzero, then $\operatorname{diagL}(T)$ is
  finite-rank, hence compact.
\end{theorem}

\begin{theorem}[Compactness criterion for block-diagonal operators]
  \label{thm:isCompactOperator-diagL}
  \lean{lpDiag.isCompactOperator_diagL}
  \leanok
  \uses{thm:isCompactOperator-diagL-of-support-finite}
  If $\|T_i\| \to 0$ along the cofinite filter on $\alpha$, then $\operatorname{diagL}(T)$ is
  a compact operator. The finitely-supported truncations of
  Theorem~\ref{thm:isCompactOperator-diagL-of-support-finite} converge to
  $\operatorname{diagL}(T)$ in operator norm.
\end{theorem}

\section{The Dirac triple of the circle \texorpdfstring{$S^1$}{S1}}
\label{sec:circle}

\textbf{Status: complete.} A fully formalized odd spectral-triple core with compact resolvent
on the circle. The algebra representation is generated by a single Fourier shift unitary; since the
Dirac eigenvalues are additive in the shift, the commutator of the shift with $D$ collapses to
a scalar multiple of the shift itself, simplifying the boundedness argument relative to the
torus case (Section~\ref{sec:torus}), which uses the same diagonal-operator technology but
with matrix-valued blocks.

The model is $H = \ell^2(\mathbb{Z})$, the Fourier side of $L^2(S^1)$, with Dirac operator
$D = -i\,d/d\theta$ acting diagonally on the orthonormal Fourier basis $(e_n)_{n \in
\mathbb{Z}}$ by $D e_n = n \cdot e_n$, with maximal domain the Sobolev space $H^1 = \{a \in
\ell^2(\mathbb{Z}) : \sum n^2 |a_n|^2 < \infty\}$.

\begin{definition}[Circle Dirac operator]
  \label{def:circle-diracDirac}
  \lean{SpectralTriples.Circle.diracDirac}
  \leanok
  The operator $D \colon H^1 \to \ell^2(\mathbb{Z})$, $(D a)_n = n \cdot a_n$.
\end{definition}

\begin{theorem}[Self-adjointness]
  \label{thm:circle-diracDirac-isSelfAdjoint}
  \lean{SpectralTriples.Circle.diracDirac_isSelfAdjoint}
  \leanok
  \uses{def:circle-diracDirac}
  The circle Dirac operator is self-adjoint.
\end{theorem}

\begin{theorem}[$i$ lies in the resolvent set]
  \label{thm:circle-mem-resolventSet-I}
  \lean{SpectralTriples.Circle.mem_resolventSet_I}
  \leanok
  \uses{thm:circle-diracDirac-isSelfAdjoint, thm:basic-criterion-self-adjoint}
  $i \in \rho(D)$, by the basic criterion of self-adjointness
  (Theorem~\ref{thm:basic-criterion-self-adjoint}) applied to
  Theorem~\ref{thm:circle-diracDirac-isSelfAdjoint}.
\end{theorem}

\begin{theorem}[Compact resolvent]
  \label{thm:circle-isCompactOperator-resolvent-I}
  \lean{SpectralTriples.Circle.isCompactOperator_resolvent_I}
  \leanok
  \uses{thm:circle-mem-resolventSet-I, thm:isCompactOperator-diagL}
  The resolvent $R(i, D)$ is a compact operator on $\ell^2(\mathbb{Z})$. It is the diagonal
  operator $b_n \mapsto b_n / (i - n)$, with block norms $1/|i-n| \to 0$, via
  Theorem~\ref{thm:isCompactOperator-diagL}.
\end{theorem}

\begin{definition}[Circle algebra and representation]
  \label{def:circle-algebra}
  \lean{SpectralTriples.Circle.algebra, SpectralTriples.Circle.rep}
  \leanok
  The star-subalgebra of $\mathcal{B}(\ell^2(\mathbb{Z}))$ generated by the single Fourier
  shift unitary $W_1$ (the Fourier image of the coordinate function $e^{i\theta}$ on $S^1$),
  represented by inclusion into $\mathcal{B}(\ell^2(\mathbb{Z}))$.
\end{definition}

\begin{theorem}[Assembly: odd spectral-triple core]
  \label{thm:circle-isOddSpectralTriple}
  \lean{SpectralTriples.Circle.isOddSpectralTriple}
  \leanok
  \uses{def:odd-spectral-triple, thm:circle-diracDirac-isSelfAdjoint, def:circle-algebra}
  The circle data $(A, \ell^2(\mathbb{Z}), \pi, D)$, with $A$ the algebra of
  Definition~\ref{def:circle-algebra}, satisfies the odd core predicate. Since the Dirac
  eigenvalues are additive in the shift index, the commutator $[D, W_g]$ collapses to the
  scalar multiple $g \cdot W_g$ of the shift itself; this bound is propagated through the generated
  star-subalgebra by \texttt{StarAlgebra.adjoin\_induction}.
\end{theorem}

\begin{theorem}[Compact resolvent]
  \label{thm:circle-isCompactResolventSpectralTriple}
  \lean{SpectralTriples.Circle.isCompactResolventSpectralTriple}
  \leanok
  \uses{thm:circle-isOddSpectralTriple, thm:circle-isCompactOperator-resolvent-I,
        def:compact-resolvent-spectral-triple}
  The circle core satisfies the compact-resolvent predicate at $i$.
\end{theorem}

\section{The Dirac triple of the two-torus \texorpdfstring{$T^2$}{T2}}
\label{sec:torus}

\textbf{Status: complete.} This is the flagship example: a fully formalized even
spectral-triple core with compact resolvent on the flat torus
$T^2 = \mathbb{R}^2/\mathbb{Z}^2$, whose graded-kernel invariant is computed and shown to
vanish.

The Hilbert space is $H = \ell^2(\mathbb{Z}^2; \mathbb{C}^2)$ (square-summable spinor-valued
sequences over the Fourier lattice $\mathbb{Z}^2$). On the Fourier mode $(m, n)$ the Dirac
operator acts on the spinor fibre $\mathbb{C}^2$ by the self-adjoint block
\[
  D_{(m,n)} = 2\pi (\sigma_1 m + \sigma_2 n)
    = 2\pi \begin{pmatrix} 0 & m - in \\ m + in & 0 \end{pmatrix},
\]
with eigenvalues $\pm 2\pi\sqrt{m^2+n^2}$; the chirality grading is $\gamma = \sigma_3 =
\operatorname{diag}(1,-1)$ fibrewise.

\begin{definition}[Torus Dirac operator]
  \label{def:torus-diracDirac}
  \lean{SpectralTriples.Torus.diracDirac}
  \leanok
  The block-diagonal operator $D \colon H^1 \to H$, $(D a)_{(m,n)} = D_{(m,n)} (a_{(m,n)})$,
  on the Sobolev domain $H^1$.
\end{definition}

\begin{theorem}[Self-adjointness]
  \label{thm:torus-diracDirac-isSelfAdjoint}
  \lean{SpectralTriples.Torus.diracDirac_isSelfAdjoint}
  \leanok
  \uses{def:torus-diracDirac}
  The torus Dirac operator is self-adjoint.
\end{theorem}

\begin{theorem}[Compact resolvent]
  \label{thm:torus-isCompactOperator-resolvent-I}
  \lean{SpectralTriples.Torus.isCompactOperator_resolvent_I}
  \leanok
  \uses{thm:torus-diracDirac-isSelfAdjoint, thm:isCompactOperator-diagL}
  The resolvent $R(i, D)$ is a compact operator on $H$. It is exhibited explicitly as a
  block-diagonal operator with block norms tending to $0$ (since $|D_{(m,n)}| \to \infty$),
  via Theorem~\ref{thm:isCompactOperator-diagL}.
\end{theorem}

\begin{definition}[Torus grading]
  \label{def:torus-grading}
  \lean{SpectralTriples.Torus.grading}
  \leanok
  The block-diagonal operator $\gamma$ acting fibrewise by $\sigma_3 =
  \operatorname{diag}(1,-1)$.
\end{definition}

\begin{lemma}[Properties of the grading]
  \label{lem:torus-grading-properties}
  \lean{SpectralTriples.Torus.isSelfAdjoint_grading, SpectralTriples.Torus.grading_mul_self,
        SpectralTriples.Torus.grading_anticomm}
  \leanok
  \uses{def:torus-grading, thm:torus-diracDirac-isSelfAdjoint}
  $\gamma$ is self-adjoint, $\gamma^2 = 1$, and $D\gamma\xi = -\gamma D\xi$ for $\xi \in \dom
  D$: each Dirac block is off-diagonal and $\sigma_3$ is diagonal, so they anticommute
  fibrewise.
\end{lemma}

\begin{definition}[Torus algebra and representation]
  \label{def:torus-algebra}
  \lean{SpectralTriples.Torus.algebra, SpectralTriples.Torus.rep}
  \leanok
  The star-subalgebra of $\mathcal{B}(H)$ generated by the two coordinate shift unitaries
  $W_{(1,0)}, W_{(0,1)}$ (the Fourier image of the trigonometric polynomials
  $\mathbb{C}[\mathbb{Z}^2]$ on $T^2$), represented by inclusion into $\mathcal{B}(H)$.
\end{definition}

\begin{theorem}[Assembly: odd and even spectral-triple cores]
  \label{thm:torus-isEvenSpectralTriple}
  \lean{SpectralTriples.Torus.isOddSpectralTriple, SpectralTriples.Torus.isEvenSpectralTriple}
  \leanok
  \uses{def:odd-spectral-triple, def:even-spectral-triple,
        thm:torus-diracDirac-isSelfAdjoint, lem:torus-grading-properties, def:torus-algebra}
  The torus data $(A, H, \pi, D)$, with $A$ the algebra of
  Definition~\ref{def:torus-algebra}, satisfies the odd core predicate; together with the
  grading $\gamma$ of Definition~\ref{def:torus-grading}, it satisfies the even core predicate. The
  bounded-commutator condition is checked by propagating a uniform commutator bound for the
  two generators $W_{(1,0)}, W_{(0,1)}$ through the generated star-subalgebra.
\end{theorem}

\begin{theorem}[Compact resolvent]
  \label{thm:torus-isCompactResolventSpectralTriple}
  \lean{SpectralTriples.Torus.isCompactResolventSpectralTriple}
  \leanok
  \uses{thm:torus-isEvenSpectralTriple, thm:torus-isCompactOperator-resolvent-I,
        def:compact-resolvent-spectral-triple}
  The torus core satisfies the compact-resolvent predicate at $i$.
\end{theorem}

\begin{theorem}[The graded-kernel invariant vanishes]
  \label{thm:torus-index-eq-zero}
  \lean{SpectralTriples.Torus.gradedKernelIndex_eq_zero}
  \leanok
  \uses{thm:torus-isEvenSpectralTriple, def:graded-kernel-index}
  $\operatorname{gkIndex}(D, \gamma) = 0$ for the flat torus Dirac core: the graded kernel of
  $D$ is exactly the zero Fourier mode, with $(\ker D)^{+}$ and $(\ker D)^{-}$ each spanned by
  one of the two spinor basis vectors at that mode, so $\dim(\ker D)^{+} = \dim(\ker D)^{-} =
  1$. This matches the Atiyah--Singer prediction for the trivial (degree-$0$) line bundle on
  $T^2$.  This theorem does not construct the chiral operator or identify the invariant with
  a Fredholm index.
\end{theorem}

\section{The unilateral shift on \texorpdfstring{$\ell^2(\mathbb{N})$}{l2(N)}}
\label{sec:shift}

\textbf{Status: complete.} A direct illustration of the Fredholm index machinery
(Section~\ref{sec:fredholm}) on a single bounded operator, independent of the spectral-triple
framework: the classical forward unilateral shift on $\ell^2(\mathbb{N})$, with Fredholm index
$-1$.

\begin{definition}[Unilateral shift]
  \label{def:shift}
  \lean{SpectralTriples.Shift.shift}
  \leanok
  The forward unilateral shift $S$ on $\ell^2(\mathbb{N})$: $(Sx)_n = x_{n-1}$ for $n \geq 1$
  and $(Sx)_0 = 0$.
\end{definition}

\begin{theorem}[The shift is Fredholm]
  \label{thm:isFredholm-shift}
  \lean{SpectralTriples.Shift.isFredholm_shift}
  \leanok
  \uses{def:shift, def:isFredholm}
  The shift has trivial kernel, closed range, and one-dimensional cokernel, and hence is
  Fredholm.
\end{theorem}

\begin{theorem}[The shift has index $-1$]
  \label{thm:fredholmIndex-shift}
  \lean{SpectralTriples.Shift.fredholmIndex_shift}
  \leanok
  \uses{thm:isFredholm-shift, def:fredholm-index}
  $\operatorname{ind}(S) = -1$. The shift is injective ($\ker S = 0$) and its range is closed
  with one-dimensional orthogonal complement spanned by the zeroth basis vector $e_0$, so
  $\operatorname{ind}(S) = \dim(\ker S) - \dim(\operatorname{coker} S) = 0 - 1 = -1$.
\end{theorem}

\section{The flux-\texorpdfstring{$k$}{k} magnetic phase model}
\label{sec:magnetic-dirac}

\textbf{Status: complete (bounded phase model).}  The formalized operator is the bounded
backward shift on $\ell^2(\mathbb N)\otimes\mathbb C^k$.  It is the expected phase of the
unbounded, $\sqrt{n+1}$-weighted Landau lowering operator, not that geometric Dirac operator
itself.  It is Fredholm of index $k$, and its kernel carries the expected finite Weyl action.
Constructing the weighted operator and the geometric line-bundle realization is deferred.

\begin{definition}[Flux-$k$ bounded magnetic phase]
  \label{def:magneticDirac}
  \lean{SpectralTriples.MagneticDirac.magneticPhase}
  \leanok
  \uses{def:shift}
  The backward shift on the Landau-level index $\mathbb{N}$, tensored with the identity on the
  $k$-dimensional guiding-center factor $\mathbb{C}^k$: a bounded operator $F_k$ on
  $\ell^2(\mathbb{N}) \otimes \mathbb{C}^k$.
\end{definition}

\begin{theorem}[The kernel is the lowest Landau level]
  \label{thm:magneticDirac-ker-finrank}
  \lean{SpectralTriples.MagneticDirac.magneticPhase_ker_finrank}
  \leanok
  \uses{def:magneticDirac}
  $\dim \ker(F_k) = k$: the kernel is exactly the lowest Landau level, identified with
  $\mathbb{C}^k$.
\end{theorem}

\begin{theorem}[The magnetic phase is Fredholm]
  \label{thm:isFredholm-magneticPhase}
  \lean{SpectralTriples.MagneticDirac.isFredholm_magneticPhase}
  \leanok
  \uses{thm:magneticDirac-ker-finrank, def:isFredholm}
  The bounded phase $F_k$ is Fredholm: its kernel is $k$-dimensional and it is surjective.
\end{theorem}

\begin{theorem}[The magnetic phase has index $k$]
  \label{thm:fredholmIndex-magneticDirac}
  \lean{SpectralTriples.MagneticDirac.fredholmIndex_magneticPhase}
  \leanok
  \uses{thm:isFredholm-magneticPhase, def:fredholm-index}
  $\operatorname{ind}(F_k) = k$: the operator is surjective (trivial cokernel), so the index
  equals the kernel dimension.
\end{theorem}

\begin{definition}[Magnetic translations]
  \label{def:magnetic-translations}
  \lean{SpectralTriples.MagneticDirac.magClock, SpectralTriples.MagneticDirac.magShift}
  \leanok
  The magnetic clock $\hat C$ and cyclic shift $\hat S$ operators on the guiding-center factor
  $\mathbb{C}^k$, acting by the phase $\omega_k = e^{2\pi i / k}$ and by cyclic permutation
  respectively.
\end{definition}

\begin{theorem}[Finite Weyl relation and commutation with the phase]
  \label{thm:magnetic-translation-weyl}
  \lean{SpectralTriples.MagneticDirac.magneticTranslation_weyl,
        SpectralTriples.MagneticDirac.magClock_comm_dirac,
        SpectralTriples.MagneticDirac.magShift_comm_dirac}
  \leanok
  \uses{def:magnetic-translations, def:magneticDirac}
  The magnetic clock and shift satisfy the finite Weyl relation $\hat C \hat S = \omega_k \hat
  S \hat C$, and both commute with $F_k$ — certifying that the $k$-dimensional kernel of
  Theorem~\ref{thm:magneticDirac-ker-finrank} carries the genuine flux-$k$ Heisenberg
  degeneracy.
\end{theorem}

\section{Theta sections on the square torus}
\label{sec:theta-sections}

\textbf{Status: complete (function-theoretic lower bound).}  The code constructs $k$ explicit
entire functions satisfying the degree-$k$ automorphy relations on the square torus and proves
them linearly independent.  It does not yet define a geometric $L^2$ line-bundle Dirac
operator or identify these functions with its kernel.  The exact dimension of the automorphic
holomorphic-function space is completed in Section~\ref{sec:fourier-holomorphic}; the operator
interpretation is deferred to the planned bridge below.

\begin{definition}[Explicit degree-$k$ theta sections]
  \label{def:thetaSection}
  \lean{SpectralTriples.ThetaSections.thetaSection}
  \leanok
  For $k \geq 1$ and $a \in \{0, \dots, k-1\}$, the theta section
  \[
    \theta_{k,a}(z) = e^{2\pi i a z} \, \vartheta(kz + ai;\, ki),
  \]
  built from Mathlib's two-variable Jacobi theta function $\vartheta$ at modulus $\tau = ki$.
\end{definition}

\begin{theorem}[Theta sections are holomorphic]
  \label{thm:differentiable-thetaSection}
  \lean{SpectralTriples.ThetaSections.differentiable_thetaSection}
  \leanok
  \uses{def:thetaSection}
  Each $\theta_{k,a}$ is differentiable (holomorphic) on $\mathbb{C}$.
\end{theorem}

\begin{theorem}[Automorphy under the square-torus lattice]
  \label{thm:thetaSection-quasiPeriodic}
  \lean{SpectralTriples.ThetaSections.thetaSection_quasiPeriodic,
        SpectralTriples.ThetaSections.thetaSection_periodic}
  \leanok
  \uses{def:thetaSection}
  $\theta_{k,a}(z+1) = \theta_{k,a}(z)$ (periodicity under the lattice generator $1$) and
  $\theta_{k,a}(z+i) = e^{-\pi i k (2z+i)} \, \theta_{k,a}(z)$ (the degree-$k$ automorphy factor
  under the generator $i$): together, these are the defining quasi-periodicity conditions for a
  holomorphic section of the degree-$k$ line bundle on $\mathbb{C}/(\mathbb{Z}+i\mathbb{Z})$.
\end{theorem}

\begin{theorem}[Diagonalization under translation by $1/k$]
  \label{thm:thetaSection-translate}
  \lean{SpectralTriples.ThetaSections.thetaSection_translate}
  \leanok
  \uses{def:thetaSection}
  $\theta_{k,a}(z + 1/k) = \omega_k^{a} \, \theta_{k,a}(z)$, where $\omega_k = e^{2\pi i / k}$:
  translation by $1/k$ diagonalizes the theta sections, with the same $k$-th roots of unity
  that appear in the magnetic clock operator of Definition~\ref{def:magnetic-translations}.
\end{theorem}

\begin{theorem}[Linear independence of the theta sections]
  \label{thm:thetaSection-linearIndependent}
  \lean{SpectralTriples.ThetaSections.thetaSection_linearIndependent}
  \leanok
  \uses{thm:thetaSection-translate}
  The $k$ automorphic entire functions $\theta_{k,0}, \dots, \theta_{k,k-1}$ are linearly
  independent over $\mathbb{C}$: they are eigenvectors of translation by $1/k$ with pairwise
  distinct eigenvalues $\omega_k^{0}, \dots, \omega_k^{k-1}$ (the $k$ distinct $k$-th roots of
  unity),
  hence independent.  After bundling them into \texttt{holSection}, this supplies the
  function-theoretic lower bound $\dim H^0(L_k)\geq k$; no operator kernel is involved.
\end{theorem}

\section{The Fourier-coefficient recursion: an exact dimension count}
\label{sec:fourier-holomorphic}

\textbf{Status: complete (function-theoretic M3a and M3b).} This section completes the
dimension count begun in Section~\ref{sec:theta-sections} (M2) by an entirely algebraic route,
with \emph{no index theorem and no $L^2$ analysis}: the upper bound
$\dim H^0(L_k) \leq k$ (M3a), an explicit theta-section basis and coordinate equivalence
$H^0(L_k)\simeq\mathbb C^k$, and the vanishing $H^0(L_{-k})=0$ (M3b).  Classical Serre
duality identifies the latter with $H^1(L_k)=0$, but that duality and the operator-cokernel
interpretation are not formalized.  These formal results rest on the contour shift,
Fourier-coefficient recursion, and the explicit theta family. See
\texttt{SpectralTriples/docs/INDEX\_PAIRING.md} for the remaining $L^2$ operator bridge.

\begin{theorem}[Contour shift]
  \label{thm:periodIntegral-eq-of-periodic}
  \lean{SpectralTriples.periodIntegral_eq_of_periodic}
  \leanok
  For an entire $1$-periodic function $f$, the period integral $\int_0^1 f(x+iy)\,dx$ does not
  depend on the height $y$. (Cauchy--Goursat on the rectangle $[0,1] \times [y_1,y_2]$: the two
  vertical sides cancel by periodicity, leaving the two horizontal integrals equal.)
\end{theorem}

\begin{definition}[Holomorphic sections of the degree-$k$ line bundle]
  \label{def:holSection}
  \lean{SpectralTriples.holSection}
  \leanok
  $H^0(L_k)$: the space of entire functions $f$ with $f(z+1) = f(z)$ and the degree-$k$
  automorphy factor $f(z+i) = e^{-\pi i k(2z+i)} f(z)$ — the same quasi-periodicity conditions
  satisfied by the theta sections of Definition~\ref{def:thetaSection}.
\end{definition}

\begin{definition}[Fourier coefficient of a periodic holomorphic function]
  \label{def:holCoeff}
  \lean{SpectralTriples.holCoeff}
  \leanok
  For a $1$-periodic function $f$, the Fourier coefficient $a_m = \int_0^1 f(x)\,
  e^{-2\pi i m x}\,dx$, computed on the real period.
\end{definition}

\begin{lemma}[The Fourier-coefficient recursion]
  \label{lem:holCoeff-recursion}
  \lean{SpectralTriples.holCoeff_recursion}
  \leanok
  \uses{def:holSection, def:holCoeff, thm:periodIntegral-eq-of-periodic}
  For $f \in H^0(L_k)$, the Fourier coefficients satisfy
  \[
    a_{m+k} = e^{-\pi(2m+k)} \, a_m.
  \]
  Comparing Fourier coefficients of the degree-$k$ automorphy relation, via the contour shift
  (Theorem~\ref{thm:periodIntegral-eq-of-periodic}) relating $f$'s Fourier coefficients on
  different horizontal lines, turns the quasi-periodicity condition into this purely algebraic
  recursion: the whole coefficient sequence is determined by $(a_0, \dots, a_{k-1})$.
\end{lemma}

\begin{theorem}[Fourier completeness: vanishing coefficients imply $f = 0$]
  \label{thm:eq-zero-of-holCoeff-eq-zero}
  \lean{SpectralTriples.eq_zero_of_holCoeff_eq_zero}
  \leanok
  \uses{def:holCoeff}
  If $f$ is entire and $1$-periodic with all Fourier coefficients $a_m = 0$, then $f = 0$:
  lifting $f$ to the circle, Mathlib's Fourier completeness gives $f$ vanishes on $\mathbb{R}$,
  and the identity theorem for the entire function $f$ then gives $f \equiv 0$ on $\mathbb{C}$.
\end{theorem}

\begin{theorem}[The upper bound: $\dim H^0(L_k) \leq k$]
  \label{thm:holSection-finrank-le}
  \lean{SpectralTriples.holSection_finrank_le}
  \leanok
  \uses{lem:holCoeff-recursion, thm:eq-zero-of-holCoeff-eq-zero}
  The restriction map $f \mapsto (a_0, \dots, a_{k-1})$ is an injective linear map
  $H^0(L_k) \to \mathbb{C}^k$: by the recursion (Lemma~\ref{lem:holCoeff-recursion}),
  $a_0 = \cdots = a_{k-1} = 0$ forces every Fourier coefficient to vanish, hence $f = 0$ by
  Theorem~\ref{thm:eq-zero-of-holCoeff-eq-zero}. Injectivity gives $\dim H^0(L_k) \leq k$.
\end{theorem}

\begin{theorem}[The exact count: $\dim H^0(L_k) = k$]
  \label{thm:holSection-finrank-eq}
  \lean{SpectralTriples.holSection_finrank_eq}
  \leanok
  \uses{thm:holSection-finrank-le, thm:thetaSection-linearIndependent}
  Combining the upper bound of Theorem~\ref{thm:holSection-finrank-le} (M3a) with the lower
  bound of Theorem~\ref{thm:thetaSection-linearIndependent} (M2, the theta sections span a
  $k$-dimensional subspace) gives $\dim H^0(L_k) = k$ exactly — a complete dimension theorem
  with no index theorem and no $L^2$ analysis.
\end{theorem}

\begin{theorem}[The theta sections form a basis]
  \label{thm:thetaHolSectionBasis}
  \lean{SpectralTriples.thetaHolSectionBasis,
        SpectralTriples.thetaHolSectionBasis_apply,
        SpectralTriples.thetaHolSectionEquiv}
  \leanok
  \uses{thm:holSection-finrank-eq, thm:thetaSection-linearIndependent}
  For $k>0$, the $k$ explicit theta sections form a basis of $H^0(L_k)$, and evaluation of
  that basis at $a\in\operatorname{Fin}(k)$ recovers the bundled theta section
  $\theta_{k,a}$.  Its coordinate map gives a linear equivalence
  $H^0(L_k)\simeq (\operatorname{Fin}(k)\to\mathbb C)$.  This completes the explicit
  function-space description; it does not yet equip the sections with the geometric $L^2$
  structure or realize them as a Dirac kernel.
\end{theorem}

\begin{definition}[Holomorphic sections of the degree-$(-k)$ line bundle]
  \label{def:holSectionNeg}
  \lean{SpectralTriples.holSectionNeg}
  \leanok
  $H^0(L_{-k})$: the same function-space definition as $H^0(L_k)$
  (Definition~\ref{def:holSection}) but with the opposite-sign automorphy factor
  $f(z+i) = e^{\pi i k(2z+i)} f(z)$.  Its classical identification with $H^1(L_k)$ uses
  Serre duality, which is not part of the Lean statement.
\end{definition}

\begin{lemma}[Parseval decay of Fourier coefficients]
  \label{lem:holCoeff-tendsto-atTop-zero}
  \lean{SpectralTriples.holCoeff_tendsto_atTop_zero}
  \leanok
  \uses{def:holCoeff}
  For any entire $1$-periodic $f$, the Fourier coefficients $a_m \to 0$ as
  $m \to +\infty$: a consequence of $\ell^2$-summability of $|a_m|^2$ (Parseval, via the
  $L^2$ lift to the circle). This is the direction asserted by the Lean theorem.
\end{lemma}

\begin{theorem}[Negative-degree holomorphic sections vanish]
  \label{thm:holSectionNeg-eq-bot}
  \lean{SpectralTriples.holSectionNeg_eq_bot, SpectralTriples.holSectionNeg_finrank_eq_zero}
  \leanok
  \uses{def:holSectionNeg, lem:holCoeff-recursion, lem:holCoeff-tendsto-atTop-zero,
        thm:eq-zero-of-holCoeff-eq-zero}
  $H^0(L_{-k}) = 0$ for $k > 0$. The opposite-sign automorphy gives the recursion
  $a_{m+k} = e^{\pi(2m+k)} a_m$,
  whose factors eventually force growth along every forward $k$-step residue progression;
  this clashes with the Parseval decay of
  Lemma~\ref{lem:holCoeff-tendsto-atTop-zero} unless every coefficient is already $0$, and then
  Theorem~\ref{thm:eq-zero-of-holCoeff-eq-zero} gives $f = 0$.  No statement about a geometric
  Dirac cokernel is formalized here.
\end{theorem}

\section{The Hermite functions: Gaussian-weighted orthogonality}
\label{sec:hermite-l2}

\textbf{Status: partial (orthonormality formalized; completeness open).}  For the
probabilists' Hermite polynomials, the normalized $L^2$ functions are
\[
  h_n(x)=\bigl(n!\sqrt{2\pi}\bigr)^{-1/2}H_n(x)e^{-x^2/4}.
\]
The exponent is $-x^2/4$, because squaring $h_n$ produces the $e^{-x^2/2}$ weight in the
polynomial orthogonality theorem.  The code now constructs these functions in $L^2(\mathbb R)$
and proves they are orthonormal.  Proving their span is dense, and hence obtaining a
\texttt{HilbertBasis}, remains the next analytic node.

\begin{lemma}[Hermite derivative identity]
  \label{lem:derivative-hermite}
  \lean{Polynomial.derivative_hermite}
  \leanok
  The probabilists' Hermite polynomials satisfy $H_{n+1}' = (n+1) \cdot H_n$. (Mathlib has the
  three-term recursion $H_{n+1} = X H_n - n H_{n-1}$ but not this derivative form.)
\end{lemma}

\begin{lemma}[Integrability against the Gaussian weight]
  \label{lem:integrable-aeval-mul-gaussian}
  \lean{Polynomial.integrable_aeval_mul_gaussian}
  \leanok
  Any polynomial $p$ times the Gaussian weight $e^{-x^2/2}$ is integrable on $\mathbb{R}$.
\end{lemma}

\begin{theorem}[Off-diagonal orthogonality]
  \label{thm:hermite-integral-eq-zero-of-ne}
  \lean{Polynomial.hermite_integral_eq_zero_of_ne}
  \leanok
  \uses{lem:derivative-hermite, lem:integrable-aeval-mul-gaussian}
  For $m \neq n$, $\int_{\mathbb{R}} H_m(x) H_n(x) e^{-x^2/2}\,dx = 0$. From the Rodrigues
  identity $(H_n \cdot w)' = -H_{n+1} \cdot w$ (with $w = e^{-x^2/2}$) and a single integration
  by parts, the weighted pairing of a polynomial $P$ against $H_{n+1}$ satisfies
  $\langle P, H_{n+1}\rangle = \langle P', H_n\rangle$; iterating shows $\deg P \leq n$ forces
  $\langle P, H_{n+1}\rangle = 0$, which gives off-diagonal vanishing without $n$-fold
  integration by parts.
\end{theorem}

\begin{theorem}[Diagonal value]
  \label{thm:hermite-integral-self}
  \lean{Polynomial.hermite_integral_self}
  \leanok
  \uses{lem:derivative-hermite, thm:hermite-integral-eq-zero-of-ne}
  $\int_{\mathbb{R}} H_n(x)^2 e^{-x^2/2}\,dx = n! \sqrt{2\pi}$. The same integration-by-parts
  recursion gives $\langle H_{n+1}, H_{n+1}\rangle = (n+1) \langle H_n, H_n \rangle$
  (Lemma~\ref{lem:derivative-hermite}), reducing to the base case
  $\int_{\mathbb{R}} e^{-x^2/2}\,dx = \sqrt{2\pi}$ (the Gaussian integral).
\end{theorem}

\begin{theorem}[Gaussian-weighted orthogonality]
  \label{thm:hermite-orthogonality}
  \lean{Polynomial.hermite_orthogonality}
  \leanok
  \uses{thm:hermite-integral-eq-zero-of-ne, thm:hermite-integral-self}
  $\int_{\mathbb{R}} H_m(x) H_n(x) e^{-x^2/2}\,dx = n! \sqrt{2\pi} \cdot \delta_{mn}$, combining
  the off-diagonal and diagonal cases.
\end{theorem}

\begin{definition}[Normalized probabilists' Hermite function]
  \label{def:hermiteFunction}
  \lean{Polynomial.hermiteFunction}
  \leanok
  \uses{thm:hermite-orthogonality}
  The normalized function is
  \[
    h_n(x)=\bigl(n!\sqrt{2\pi}\bigr)^{-1/2}H_n(x)e^{-x^2/4}.
  \]
\end{definition}

\begin{definition}[Hermite function as an $L^2$ vector]
  \label{def:hermiteFunctionL2}
  \lean{Polynomial.hermiteFunctionL2}
  \leanok
  \uses{def:hermiteFunction, lem:integrable-aeval-mul-gaussian}
  Each $h_n$ is square-integrable and therefore defines a vector in scalar-valued
  $L^2(\mathbb R)$ over either $\mathbb R$ or $\mathbb C$.
\end{definition}

\begin{theorem}[The normalized Hermite functions are orthonormal]
  \label{thm:orthonormal-hermiteFunctionL2}
  \lean{Polynomial.orthonormal_hermiteFunctionL2}
  \leanok
  \uses{def:hermiteFunctionL2, thm:hermite-orthogonality}
  The family $(h_n)_{n\in\mathbb N}$ is orthonormal in $L^2(\mathbb R)$.
\end{theorem}

\begin{theorem}[Completeness of the Hermite functions (planned)]
  \label{thm:hermite-completeness}
  \uses{thm:orthonormal-hermiteFunctionL2}
  The closed linear span of $(h_n)_{n\in\mathbb N}$ is all of $L^2(\mathbb R)$, so the
  orthonormal family packages as a \texttt{HilbertBasis}.  This density statement is not yet
  formalized.
\end{theorem}

\section{Planned geometric operator bridge}
\label{sec:planned-operator-bridge}

The basis of Theorem~\ref{thm:thetaHolSectionBasis} completes the holomorphic
automorphic-function calculation.  What remains is analytic rather than a further dimension
count: construct the geometric $L^2$ spaces and closed chiral operator, identify the theta
basis with its zero modes, and connect that operator to the Landau-level model.

\begin{definition}[The $L^2$ line-bundle Dirac operator (planned)]
  \label{def:line-bundle-dirac-L2}
  \uses{def:holSection, thm:thetaHolSectionBasis, def:holSectionNeg,
        def:compact-resolvent-spectral-triple}
  Construct the weighted Hilbert spaces $L^2(T^2;S^\pm\otimes L_k)$ of quasi-periodic
  sections and the closed, densely defined twisted chiral operator
  $D^+_{L_k}\colon\dom(D^+_{L_k})\to L^2(T^2;S^-\otimes L_k)$.  No such space or operator is
  currently present in Lean.
\end{definition}

\begin{theorem}[Weighted Landau lowering bridge (planned)]
  \label{thm:weighted-landau-lowering}
  \uses{def:line-bundle-dirac-L2, thm:hermite-completeness, def:magneticDirac,
        thm:isFredholm-magneticPhase}
  Diagonalize $D^+_{L_k}$ as the unbounded weighted lowering operator
  $a_ke_{n+1,j}=\sqrt{n+1}\,e_{n,j}$.  Prove that its phase is the bounded backward shift
  $F_k$ of Definition~\ref{def:magneticDirac}; the unbounded operator $D^+_{L_k}$ itself is not
  unitarily equivalent to $F_k$.  This bridge will transport the already-formalized kernel,
  cokernel, and index of the phase model to the geometric operator.
\end{theorem}
